\label{chapter:Experiment}

To demonstrate the relatively concrete performance of our protocol, we have developed a prototype
implementation in C++. The source code is available at~\url{[TODO: insert repository URL]}.
Compared to other protocols relying on pairing-based assumptions, our lattice-based construction
cannot match their concrete efficiency; pairing groups offer compact group elements and fast
bilinear-map evaluations that are simply not available in the lattice setting. However, the proposed
scheme provides post-quantum security---a property that no pairing-based threshold batch IBE can
offer---together with fully non-interactive pre-decryption, and the prototype shows that even at a
toy parameter scale the construction is feasible and its bottlenecks are well-localised.

This chapter is organised as follows. Section~\ref{sec:asymptotic} performs an asymptotic and
concrete size analysis, compares the output sizes against the threshold batch IBE of Boneh, Laufer,
and Tas~\cite{EPRINT:BonLauTas25} (\textsf{BLT25}), and discusses the structural trade-offs between
the two thresholdization frameworks. Section~\ref{sec:implementation} describes the implementation,
the experimental setup, and the measured results. Table~\ref{tab:eval-metrics} summarises the
evaluation metrics measured on the prototype; Section~\ref{subsec:secure-params} presents the
proposed production-grade parameter set targeting $2^{64}$ security and analyses its hardness
using the lattice estimator of~\cite{EC:ADPS16}.

\begin{table}[h]
\centering
\renewcommand{\arraystretch}{1.25}
\begin{tabular}{|l|l|r|}
\hline
\textbf{Metric} & \textbf{Unit} & \textbf{Measured value} \\
\hline
Public-key size $|\mathsf{pk}|$                    & B   & 847 \\
\hline
Ciphertext size $|\mathsf{ct}|$                    & B   & 119 \\
\hline
Per-party pre-decryption key $|\mathsf{sbk}_j|$    & B   & 14 \\
\hline
Combined pre-decryption key $|\mathsf{sbk}|$       & B   & 14 \\
\hline
Per-party secret key $|\mathsf{sk}_j|$             & B   & 56 \\
\hline
$\mathsf{Setup}$ time                              & µs  & 8.87 \\
\hline
$\mathsf{Enc}$ time (per message)                  & µs  & 71.05 \\
\hline
$\mathsf{PreDec}$ time (per party)                 & µs  & 67.62 \\
\hline
\quad of which $\mathsf{SampleLeft}$               & µs  & 66.79 \\
\hline
\quad of which $\mathsf{PSampPre}$                 & µs  & 0.83 \\
\hline
$\mathsf{Combine}$ time                            & µs  & 0.39 \\
\hline
$\mathsf{Dec}$ time (per ciphertext)               & µs  & 59.35 \\
\hline
Pre-decryption communication (per party)            & B   & 14 \\
\hline
Total committee communication (per batch)           & B   & 42 \\
\hline
\end{tabular}
\caption{Evaluation metrics for the toy-PoC ($Q=12289$, $n=4$, $k=3$, $N=5$, batch
size $\ell=5$). Hash oracles backed by SHAKE-256 (Keccak-1600, FIPS~202).
Timings are wall-clock means over 500–2\,000 repetitions; sizes are bit-packed logical
sizes ($\lceil\log_2 Q\rceil=14$ bits per $\mathbb{Z}_q$ element).}
\label{tab:eval-metrics}
\end{table}

\section{Asymptotic Analysis}
\label{sec:asymptotic}

This section derives the asymptotic and concrete sizes of the scheme outputs, compares them
against \textsf{BLT25}, and analyses the structural trade-offs between the two thresholdization
approaches.

\subsection{Parameter Instantiation}
\label{subsec:param-inst}

We instantiate the scheme at a toy security level to make the dimensions and Gaussian widths
concrete. The chosen parameters are listed in Table~\ref{tab:toy-params}.

\begin{table}[h]
\centering
\renewcommand{\arraystretch}{1.25}
\begin{tabular}{|c|c|c|c|c|c|}
\hline
$\lambda$ & $n$ & $q$ & $k$ & $N$ & $\ell$ \\
\hline
$8$ & $8$ & $2^{96}-17$ & $2$ & $3$ & $2$ \\
\hline
\end{tabular}
\caption{Toy parameter instantiation.}
\label{tab:toy-params}
\end{table}

From these values and the identity encoding length $d = 1$ we derive the lattice dimensions
\[
m = m_{\mathbf G} + n = 776, \qquad
\widetilde m = nk\lceil\log q\rceil = 1536, \qquad
\overline m = 3nk\lceil\log q\rceil = 4608, \qquad
t = \widetilde m(d+1) = 3072,
\]
and the Gaussian widths $\sigma_1 \approx 4.4\times 10^{9}$ (for explainable $\mathsf{SampleLeft}$)
and $\sigma_2 \approx 6.0\times 10^{19}$ (for $\mathsf{PSampPre}$), with correctness bound
$\beta \approx 6.7\times 10^{21}$.

\subsection{Lattice Security Estimator}
\label{subsec:estimator}

\paragraph{Toolkit overview.}
We use the open-source \emph{lattice estimator}~\cite{EC:ADPS16} (repository:
\texttt{github.com/malb/lattice-estimator}) to estimate the concrete bit-security of the
LWE and SIS instances that arise in TBIBE. The tool, implemented in SageMath, collects a
catalogue of the best-known classical and quantum lattice attacks against a given set of
parameters and reports the estimated cost of the cheapest one. It has become the de-facto
standard for parameter selection in lattice-based cryptography~\cite{NIST:PQC-round3}.

\paragraph{LWE hardness model.}
Given an LWE instance
$\mathsf{LWE}_{n,q,\chi_s,\chi_e,m}$---dimension $n$, modulus $q$, secret distribution
$\chi_s$, error distribution $\chi_e$, and $m$ available samples---the estimator reduces
the problem to finding short vectors in a structured lattice derived from the samples and
then applies cost models for the BKZ lattice-reduction algorithm~\cite{C:CheNgu11}.
The central quantity is the \emph{BKZ block size} $\beta$: the size of the largest
sublattice enumerated at each BKZ tour. A larger $\beta$ means a stronger (more expensive)
reduction is needed to distinguish the LWE samples from random, and the estimated attack cost
grows roughly as $2^{0.292\beta}$ (the MATZOV bound~\cite{MATZOV22}).
The estimator evaluates several attack strategies (dual attack, primal uSVP, primal hybrid,
lattice sieving) and returns the minimum-cost strategy together with its BKZ block size
and operation count (\emph{rop}, ring operations).

\paragraph{SIS hardness model.}
For the SIS problem $\mathsf{SIS}_{n,q,\beta,m}$---find a non-zero
$\mathbf{v}\in\mathbb{Z}^m$ with $\mathbf{A}\mathbf{v}\equiv\mathbf{0}\pmod{q}$ and
$\|\mathbf{v}\|\leq\beta$---the hardness depends on whether $\beta$ is smaller than the
Gaussian heuristic for the shortest vector in the $q$-ary lattice
$\Lambda_q^\perp(\mathbf{A})$:
\[
    \lambda_1\bigl(\Lambda_q^\perp(\mathbf{A})\bigr)
    \;\approx\; \sqrt{\tfrac{m}{2\pi e}}\cdot q^{n/m}.
\]
When $\beta > \lambda_1$ the problem is trivially easy (any lattice vector is short enough);
when $\beta < \lambda_1$ BKZ must find a genuinely short vector and the cost is determined
by the same block-size analysis as LWE.

\paragraph{TBIBE LWE instance.}
In TBIBE, the ciphertext $(\mathbf{b}_1,\mathbf{b}_2,\mathbf{b}_3)$ is produced by LWE
encryption with secret $\mathbf{s}\in\mathbb{Z}_q^{nk}$, noise drawn from $\chi_0$, and
$2mk$ samples (the columns of $\mathbf{C}$). We model this as
$\mathsf{LWE}_{nk,\,q,\,\mathsf{Uniform}(\mathbb{Z}_q),\,\mathcal{D}_{\chi_0},\,2mk}$,
where the maximum noise $\chi_0$ compatible with correct decryption is
$\chi_0 \leq q\,/\,(4\beta\sqrt{2mk})$. We pass this maximum-noise instance to the
estimator, giving a conservative (lower-bound) security estimate: any smaller noise only
makes the LWE instance easier to attack.

\subsection{Proposed Secure Parameter Set}
\label{subsec:secure-params}

The toy instantiation is suitable only for validating the algebraic correctness of the
construction. To identify a parameter set that achieves a concrete security level of $2^{64}$
operations against lattice attacks, we ran a systematic search over the LWE and SIS hardness
of the underlying assumptions using the lattice estimator (Section~\ref{subsec:estimator},
MATZOV cost model, BKZ algorithm family). The search fixes $k=2$, $N=5$, $\ell=4$, and scans
$n\in\{32,64,96,\ldots,1024,1280\}$; for each $n$ the minimum modulus size
$\lceil\log_2 q\rceil$ that satisfies the correctness bound
$q > 4\beta\sqrt{2mk}$ is derived analytically before the estimator is invoked.

\paragraph{Hardness landscape.}
A key structural constraint limits security at small dimensions. The scheme's correctness
condition forces $\lceil\log_2 q\rceil \gtrsim 90$ for any $n$, which implies
$2mk / nk \approx 2\lceil\log_2 q\rceil \geq 180$. This large ratio of samples to unknowns
places the resulting LWE instance in the \emph{overdetermined} regime, where the BKZ
algorithm achieves its minimum useful block size ($\beta=40$) regardless of $n$.
Consequently, every parameter set with $n\leq 768$ yields only $2^{33}$--$2^{55}$
security against raw LWE attacks. Security crosses $2^{64}$ only at $n=1280$ (dimension
$nk=2560$), where BKZ requires block size $\beta=118$ and $2^{67.4}$ operations.
Table~\ref{tab:secure-params} lists the proposed $2^{64}$ parameter set and all derived
dimensions.

\begin{table}[h]
\centering
\renewcommand{\arraystretch}{1.25}
\begin{tabular}{|l|r|l|}
\hline
\textbf{Parameter} & \textbf{Value} & \textbf{Role} \\
\hline
$n$ & $1280$ & LWE dimension \\
\hline
$k$ & $2$ & partial-trapdoor width \\
\hline
$N,\,\ell$ & $5,\,4$ & parties, batch size \\
\hline
$\lceil\log_2 q\rceil$ & $127$ & modulus size (bits) \\
\hline
$nk$ & $2560$ & effective LWE dimension \\
\hline
$m = m_{\mathbf G}+n$ & $163\,840$ & base lattice columns \\
\hline
$2mk$ & $655\,360$ & LWE sample count \\
\hline
$\widetilde m = nk\lceil\log q\rceil$ & $325\,120$ & \\
\hline
$\overline m = 3nk\lceil\log q\rceil$ & $975\,360$ & \\
\hline
$t = 2\widetilde m$ & $650\,240$ & identity expansion length \\
\hline
$\sigma_2$ & $\approx 2^{104}$ & PSampPre Gaussian width \\
\hline
$\beta$ & $\approx 2^{114.7}$ & correctness bound on $\|\mathbf{sbk}\|$ \\
\hline
$\sigma_1$ & $\approx 2^{45.9}$ & SampleLeft Gaussian width \\
\hline
\hline
$|\mathsf{pk}|$ & $\approx 27.71\ \mathrm{GB}$ & public key \\
\hline
$|\mathsf{ct}|$ & $\approx 34.5\ \mathrm{MB}$ & ciphertext (batch $\ell=4$) \\
\hline
$|\mathsf{sbk}_j|$ & $\approx 9.0\ \mathrm{MB}$ & per-party pre-decryption key \\
\hline
\hline
\multicolumn{3}{|l|}{\textbf{Security (lattice estimator, MATZOV BKZ cost model)}} \\
\hline
LWE hardness & $2^{67.4}$ operations & BKZ $\beta=118$ \\
\hline
\end{tabular}
\caption{Proposed $2^{64}$-secure parameter set ($n=1280$, $k=2$, $\lceil\log_2 q\rceil=127$).
All derived dimensions follow from the formulas in Table~\ref{tab:parameters}.
Sizes use 127-bit elements packed without padding.}
\label{tab:secure-params}
\end{table}

\paragraph{Reduction from LWE to $\kappa$-LWE~\cite{AC:ALLW25}.}
\label{par:thm1-allw25}
The hardness of $\kappa$-LWE was first established by Ling, Phan, Stehlé, and
Steinfeld~\cite{C:LPSS14} via a PPT reduction from standard LWE, under the assumption
that all hint covariance matrices are identical ($\mathbf{S}_i = \mathbf{S}$).
The construction in this thesis uses the $\kappa$-LWE definition of Albrecht, Lai,
Lapiha, and Woo~\cite{AC:ALLW25}, whose Theorem~1 generalises the LPSS14 reduction
to \emph{per-hint} covariances $\mathbf{S}_i = \mathrm{diag}(\sigma_i\mathbf{I},\sigma_i'\mathbf{I})$.
The five conditions required by Theorem~1 are:
\begin{enumerate}[label=(A\arabic*)]
  \item $\kappa \geq 100$; all $\sigma_i,\sigma_i'>0$,
  \item $m \geq \max\!\bigl\{\Omega\!\left(\kappa\log(\kappa\max_i\sigma_i)\right),\;\Omega(n\log q)\bigr\}$,
  \item $q > \max_i(\sigma_i')\log m$\quad (prime),
  \item $\min_i(\sigma_i) > \max\!\left(\Omega\!\bigl(\sqrt{\kappa m\log m}\bigr),\;
        \max_i(\sigma_i')^{\kappa/(m+\kappa)}\right)$,
  \item $\min_i(\sigma_i') \geq \Omega\!\left(\kappa\sqrt{m}\cdot\max_i(\sigma_i)^2
        \cdot\log^{3/2}\!\bigl(\kappa m\max_i(\sigma_i)\bigr)\right)$.
\end{enumerate}
The key improvement over~\cite{C:LPSS14} is in (A5): the original coefficient $m^3$
is replaced by $\kappa\sqrt{m}$, a $2^{29.9}$-fold reduction for the parameters below.

We map TBIBE's parameters to Theorem~1's notation.
In the ALLW25 $\kappa$-LWE formulation the LWE matrix
$\mathbf{A}\in\mathbb{Z}_q^{nk\times 2mk}$ provides $2mk$ samples, and the $\kappa$
hint vectors $\mathbf{e}_i\leftarrow\mathcal{D}_{\Lambda_q^{\perp}(\mathbf{A}),\mathbf{S}_i}$
are additional kernel objects.
Hence:
\[
  n_{\mathsf{Thm}} = nk = 2560,
  \quad
  m_{\mathsf{Thm}} = 2mk = 655\,360,
  \quad
  \kappa = 2m(k-1) = 327\,680,
  \quad
  \frac{\kappa}{m_{\mathsf{Thm}}+\kappa} = \tfrac{1}{3}.
\]
The per-hint Gaussian widths $\sigma_i$ and $\sigma_i'$ correspond to
$s_{\max}(\mathbf{S}_{\mathrm{Hon}})$ and $s_{\max}(\mathbf{S}_{\mathrm{Cor}})$ for honest
and corrupt hints respectively.
From Table~\ref{tab:parameters}: $\sigma \approx 2^{13.2}$ and $\sigma' \approx 2^{54.5}$
at the production parameter set ($n=1280$, $\log q=127$).

Conditions (A1) ($\kappa=327\,680\geq100$) and (A3)
($q/(\sigma'\log m_{\mathsf{Thm}})\approx 2^{68}\gg 1$) are satisfied comfortably.
Condition (A2) holds for the $\Omega(n\log q)$ term ($m_{\mathsf{Thm}}/(nk\cdot\lceil\log_2 q\rceil)\approx 2.02$).
Conditions (A4) and (A5) are violated: (A4) requires $\sigma > 2^{20.9}$ (have $2^{13.2}$)
and (A5) requires $\sigma' \geq 2^{62.8}$ (have $2^{54.5}$).

To find the tightest $(\sigma_{\mathsf{fp}},\sigma_{\mathsf{fp}}')$ that could satisfy
(A4)+(A5) simultaneously, we solve the fixed-point system.
Condition (A4) is binding when $\sigma_{\mathsf{fp}} = (\sigma_{\mathsf{fp}}')^{1/3}$.
Substituting into (A5):
\[
  \bigl(\sigma_{\mathsf{fp}}'\bigr)^{1/3} \;=\; C_5,
  \qquad
  C_5 \;=\; \kappa\sqrt{m_{\mathsf{Thm}}}\cdot\log^{3/2}\!\bigl(\kappa\,m_{\mathsf{Thm}}\cdot C_5\bigr),
\]
which converges iteratively to $C_5 \approx 2^{37.3}$, giving
\[
  \sigma_{\mathsf{fp}} \;\approx\; 2^{37.3},
  \qquad
  \sigma_{\mathsf{fp}}' \;\approx\; 2^{112.0}.
\]
Condition (A3) is satisfied at this fixed point ($\sigma_{\mathsf{fp}}'\log m_{\mathsf{Thm}}\approx 2^{116.3}<q=2^{127}$).
However, applying $\sigma_{\mathsf{Cor}} = 2^{112}$ in the TBIBE correctness chain gives
\[
  B_U \approx 2^{129.8},\quad
  \sigma_2 \approx 2^{161.4},\quad
  \beta \approx 2^{172.1},\quad
  q_{\mathsf{corr}} = 4\beta\sqrt{2mk} \approx 2^{183.7},
\]
so correctness requires $\log q \geq 184$.

\medskip\noindent
\textbf{Cubic advantage loss.}
Beyond the parameter constraints, the reduction itself incurs a significant tightness loss.
LPSS14 Theorem~18 states concretely~\cite{C:LPSS14}: given a $(k,\mathbf{S})$-LWE adversary
with advantage $\varepsilon$, the reduction yields an LWE adversary with advantage
$\varepsilon_0 = \Omega((\varepsilon')^3)$, where $\varepsilon' = \varepsilon - 2^{-\Omega(n)}$.
Since ALLW25 Theorem~1 follows exactly the same hybrid strategy~\cite{AC:ALLW25},
it inherits the same cubic loss.
Inverted for hardness: if standard LWE is $2^\lambda$-secure (no adversary has advantage
$> 2^{-\lambda}$), then the reduction guarantees $\kappa$-LWE is at most $2^{\lambda/3}$-secure
via this path.
Concretely, to obtain $2^{64}$ security for $\kappa$-LWE through Theorem~1 one must start
from a $2^{192}$-hard LWE instance.

A systematic scan confirms that parameter sets satisfying all five Theorem~1 conditions
exist (e.g.\ $n=2048$, $\log q=200$, raw LWE security $2^{71.8}$), but after accounting
for the cubic loss, the effective $\kappa$-LWE security is only $\approx 2^{71.8/3}\approx 2^{24}$.
Table~\ref{tab:thm1-scan} summarises the joint feasibility scan with both figures.

\begin{table}[h]
\centering
\renewcommand{\arraystretch}{1.25}
\begin{tabular}{|c|c|c|c|c|c|c|}
\hline
$n$ & $\log q$ & A3--A5 & $\chi_{0,\max}$ & LWE sec. & $\kappa$-LWE sec.$^{\varepsilon^3}$ & $|\mathsf{pk}|$ \\
\hline
$128$  & $160$ & \checkmark & $2^{1.5}$  & $\approx 2^{44}$\textsuperscript{$\dagger$} & $\approx 2^{15}$ & $0.5\ \mathrm{GB}$  \\
$512$  & $184$ & \checkmark & $2^{7.1}$  & $\approx 2^{44}$\textsuperscript{$\dagger$} & $\approx 2^{15}$ & $10\ \mathrm{GB}$   \\
$2048$ & $200$ & \checkmark & $2^{5.5}$  & $2^{71.8}$  & $\approx 2^{24}$ & $188\ \mathrm{GB}$ \\
$4096$ & $208$ & \checkmark & $2^{4.7}$  & $2^{137.7}$ & $\approx 2^{46}$ & $817\ \mathrm{GB}$ \\
\hline
$1280$ & $127$ & \multicolumn{2}{c|}{(Thm~1 inapplicable)} & $2^{67.3}$ & $2^{67.3}$ (standalone) & $27.7\ \mathrm{GB}$ \\
\hline
\end{tabular}
\caption{Joint scan: ALLW25 Theorem~1 conditions (A3)--(A5) satisfied for rows 1--4.
LWE security from the lattice estimator~\cite{EC:ADPS16};
$\kappa$-LWE security column accounts for the cubic advantage loss $\varepsilon_0=\Omega(\varepsilon^3)$
from the LPSS14/ALLW25 reduction (i.e.\ $\lambda_{\kappa\text{-LWE}}\approx\lambda_{\mathsf{LWE}}/3$).
$\dagger$: BKZ block size $\beta=40$ (minimum).
Last row: production set — $\kappa$-LWE taken as a standalone assumption, no reduction loss.}
\label{tab:thm1-scan}
\end{table}

\medskip\noindent
\textbf{Implication.}
No parameter set satisfying ALLW25 Theorem~1 achieves $2^{64}$ $\kappa$-LWE security
through the reduction: the cubic advantage loss alone requires a $2^{192}$-hard LWE
baseline, which in turn demands a far larger $n$ and $\log q$, with public-key sizes
in the tens of terabytes.
The production set ($n=1280$, $|\mathsf{pk}|=27.7\ \mathrm{GB}$, security $2^{67.3}$)
instead treats $\kappa$-LWE as a \emph{standalone computational assumption}, avoiding
the cubic loss entirely.
This is a principled design choice: standalone $\kappa$-LWE has been studied in the
literature~\cite{C:LPSS14,AC:ALLW25} and is widely believed to be hard for
$\kappa=\Theta(m)$; the analogous $k$-SIS assumption is similarly taken as a
standalone assumption in multi-key lattice constructions~\cite{C:LPSS14}.
Whether a tight (or at least polynomial-loss) reduction from standard LWE to
$\kappa$-LWE exists for $\kappa=\Theta(m)$ remains an open problem.

\paragraph{Security caveats.}
Three additional caveats apply to the estimator figures above.
\begin{enumerate}
  \item \emph{$\kappa$-LWE vs.\ raw LWE.}
    The actual hardness assumption is $\kappa$-LWE~\cite{AC:ALLW25}, a structured
    assumption where the adversary receives $\kappa = 327\,680$ hint vectors from the
    partial-trapdoor lattice. The lattice estimator evaluates only unstructured LWE,
    so the $2^{67.4}$ figure is a lower bound on security against generic lattice attacks.
  \item \emph{Overdetermined regime.}
    For $n\leq 768$ the LWE instance sits in the overdetermined regime ($2mk\gg nk$),
    and the estimator assigns BKZ block size 40 (the minimum) regardless of dimension.
    This does \emph{not} mean those sets are insecure under $\kappa$-LWE; it means the
    reduction to raw LWE cannot be used to justify their security through the estimator.
  \item \emph{Benchmark infeasibility at $2^{64}$ scale.}
    The public-key matrix $\mathbf{C}\in\mathbb{Z}_q^{2560\times 655360}$ at 127 bits
    per element requires approximately $26.8\ \mathrm{GB}$ of RAM, placing any direct
    timing measurement beyond the resources of the available machine (8 GB RAM).
    The benchmark results in Table~\ref{tab:eval-metrics} are therefore provided
    exclusively for the toy prototype; Table~\ref{tab:secure-params} gives the
    analytically computed output sizes for the production parameter set.
\end{enumerate}

\subsection{Output-Size Expressions}
\label{subsec:size-expressions}

The sizes of the four main outputs are determined by the lattice dimensions and the modulus:
\[
\left\{
\begin{aligned}
|\mathsf{pk}|
&= n(2t+m)\log q, \\
|\mathsf{ct}|
&= (1+2mk+t+\overline m)\log q, \\
|\mathsf{sbk}_j|
&= 2mk \cdot \lceil\log_2\sigma_2\rceil, \\
|\mathsf{sbk}|
&= 2mk \cdot \lceil\log_2\beta\rceil.
\end{aligned}
\right.
\]
In the first expression, $n(2t+m)$ is the number of ring elements in the public key
$\mathsf{pk} = (\mathsf{crs}_{\mathsf{samp}}, C, u)$, where the CRS of the shifted multi-preimage
sampler dominates. In the second, $2mk$ is the contribution of $C^\top s + \varepsilon_2$ and
$t + \overline m$ the contribution of $F_r^\top s + \varepsilon_3$. The pre-decryption key is a vector
of dimension $2mk$: each coordinate of $\mathsf{sbk}_j$ is a discrete Gaussian sample of width
$\sigma_2$, requiring $\lceil\log_2\sigma_2\rceil$ bits; the combined key has the same dimension but
coordinates bounded by $\beta$.

Evaluating at the toy parameters of Table~\ref{tab:toy-params}:
\[
\begin{aligned}
|\mathsf{pk}|
&= 8 \times (2\times 3072 + 776)\times 96
= 5{,}314{,}560 \text{ bits}
\approx 649 \text{ KB}, \\
|\mathsf{ct}|
&= (1 + 2\times 776\times 2 + 3072 + 4608)\times 96
= 1{,}035{,}360 \text{ bits}
\approx 126 \text{ KB}, \\
|\mathsf{sbk}_j|
&= 2\times 776\times 2 \times 68
= 211{,}072 \text{ bits}
\approx 26 \text{ KB}, \\
|\mathsf{sbk}|
&= 2\times 776\times 2 \times 73
= 226{,}304 \text{ bits}
\approx 28 \text{ KB}.
\end{aligned}
\]

\subsection{Comparison with \textsf{BLT25}}
\label{subsec:comparison-blt25}

Both schemes share the same non-threshold BIBE core---the shifted multi-preimage sampler,
explainable $\mathsf{SampleLeft}$, and dual-Regev decryption---so differences arise solely
from the thresholdization layer.

\paragraph{BLT25 thresholding technique.}
In the non-threshold BIBE of~\cite{EPRINT:BonLauTas25}, a single trusted authority holds the
trapdoor $\mathbf{T}_{\mathbf{C}}$ and calls
$\mathsf{SampPre}(\mathbf{C},\mathbf{T}_{\mathbf{C}},\mathbf{c},\sigma)$
to produce the short pre-decryption key $\mathsf{sbk}\in\mathbb{Z}_q^{m}$.
To distribute this authority among $N$ parties with reconstruction threshold $k$, \textsf{BLT25}
applies \emph{Shamir secret sharing directly on the trapdoor matrix}
$\mathbf{T}_{\mathbf{C}}\in\mathbb{Z}_q^{m\times\widetilde m'}$: each party $j$ receives a Shamir
share $\mathbf{T}_{\mathbf{C},j}$, together with $\Theta(N)$ pairwise PRF seeds shared with every
other committee member. Pre-decryption then proceeds in two rounds. In the first round every party
$j$ in a quorum $T$ locally computes a partial preimage contribution $\mathbf{p}_j$ from its
trapdoor share and blinds it with a row mask $\mathbf{m}_j$ derived from the pairwise PRF
seeds; it then broadcasts the masked syndrome $\mathbf{e}_j = \mathbf{C}\mathbf{p}_j + \mathbf{m}_j$.
In the second round the target $\mathbf{c}$ is deterministically derived from a transcript tag
$H_{\mathrm{tr}}(\mathsf{sid},\mathsf{act},(\mathbf{e}_j)_{j\in T})$ binding the first-round
syndromes, after which each party broadcasts a Lagrange-weighted preimage contribution; the
combiner strips the pairwise masks and sums to recover $\mathsf{sbk}$.
The need to bind the target to the first-round transcript is what forces two rounds: the simulator
must observe all $\mathbf{e}_j$ before it can program the target, so the rounds cannot be
collapsed. Noise flooding with $\sigma_{\mathrm{flood}} = 2^{\Omega(\lambda)}$ is required to
statistically hide each party's trapdoor share within the IND-CPA reduction.

\paragraph{Our thresholding technique.}
Rather than secret-sharing the trapdoor, we replace the single $\mathsf{SampPre}$ call with
the partial-trapdoor toolkit of~\cite{AC:ALLW25}: each party $j$ holds a structurally independent
partial trapdoor $\mathsf{ptd}_j = (\mathbf{U}_j,\mathbf{C}_j,\mathbf{T}_{\mathbf{C}_j})$
satisfying $\mathbf{C}\mathbf{U}_j \equiv \mathbf{v}_j\otimes\mathbf{G}\pmod{q}$, where
$\mathbf{v}_j$ is the $j$-th Vandermonde column. The target $\mathbf{c}$ is derived
non-interactively from the batch identities via the random oracle, so all parties agree on it
without any communication. Each party independently calls $\mathsf{PSampPre}$ to obtain
$\mathbf{sbk}_j\in\mathbb{Z}_q^{2mk}$, and $\mathsf{Rec}$ reconstructs the full key by plain
summation. This removes both rounds of interaction and the pairwise PRF infrastructure, at the
cost of enlarging the lattice dimension by a factor $k$.

Table~\ref{tab:size-vs-blt} summarises the resulting asymptotic output-size expressions.

\begin{table}[h]
\centering
\renewcommand{\arraystretch}{1.3}
\begin{tabular}{|l|l|l|}
\hline
\textbf{Output} & \textbf{This work} & \textbf{\textsf{BLT25}~\cite{EPRINT:BonLauTas25}} \\
\hline
$|\mathsf{pk}|$
& $n(2t+m)\log q$, \quad $\widetilde m = nk\lceil\log q\rceil$
& $n(2t'+m)\log q$, \quad $\widetilde m' = n\lceil\log q\rceil$ \\
\hline
$|\mathsf{ct}|$
& $(1+2mk+t+\overline m)\log q$
& $(1+m+t'+\overline m')\log q$ \\
\hline
$|\mathsf{sbk}|$
& $2mk\cdot\lceil\log_2\beta\rceil$
& $m\cdot\lceil\log_2 B_{\mathrm{th}}\rceil$ \\
\hline
$|\mathsf{sk}_j|$
& $O(mk^2\log q)$, \emph{no seeds}
& $O(m\widetilde m'\log q) + \Theta(N\kappa)$ \\
\hline
\end{tabular}
\caption{Asymptotic output-size expressions. Primed quantities ($t'$, $\widetilde m'$, $\overline m'$)
denote \textsf{BLT25} dimensions; they equal ours but with $k=1$, i.e.\
$\widetilde m' = n\lceil\log q\rceil$ versus $\widetilde m = nk\lceil\log q\rceil$.}
\label{tab:size-vs-blt}
\end{table}

\paragraph{Asymptotic impact.}
Because the partial-trapdoor matrix $\mathbf{C}$ lives in
$\mathbb{Z}_q^{nk\times 2mk}$, every derived dimension is inflated by $k$:
$\widetilde m = k\widetilde m'$, $t = kt'$, $\overline m = k\overline m'$, and
$\dim(\mathsf{sbk}) = 2mk$ versus $m$. All three size expressions in
Table~\ref{tab:size-vs-blt} therefore carry a $\Theta(k)$ multiplicative overhead.
In contrast, \textsf{BLT25} keeps the lattice at its base dimension $m$,
but each party's secret key grows as $O(m\widetilde m'\log q) + \Theta(N\kappa)$
due to the Shamir share and pairwise PRF seeds, and the pre-decryption key requires
$\sigma_{\mathrm{flood}}=2^{\Omega(\lambda)}$ flooding, forcing a factor-$\lambda$
blow-up on every threshold parameter. Our construction trades the $\lambda$-factor
flooding for a $k$-factor dimension enlargement, eliminates the $\Theta(N\kappa)$ seed
term from $|\mathsf{sk}_j|$, and reduces interaction from two rounds to zero, at the
cost of the stronger $\kappa$-LWE/$\kappa$-SIS assumption and a highly selective
security model.


\section{Implementation Details}
\label{sec:implementation}

\subsection{Prototype overview}

We implemented all five algorithms of the proposed TBIBE scheme as a self-contained
C++17 prototype (\texttt{btibe\_poc}), requiring no external cryptographic libraries.
The implementation consists of six compilation units:
\begin{itemize}
    \item \texttt{mod\_arith}: modular arithmetic over $\mathbb{Z}_q$ (addition,
          multiplication, Fermat-little-theorem inversion, matrix--vector products);
    \item \texttt{thesis\_sampler}: pseudo-random sampling (uniform $\mathbb{Z}_q$,
          centered discrete Gaussian via rounded normal);
    \item \texttt{keccak}: vendored XKCP ``more-compact'' Keccak-1600 permutation
          and SHAKE-256 extendable-output function (CC0 / public domain, $\sim$30 lines);
    \item \texttt{toy\_hash}: deterministic hash functions modelling $H$, $H_{\mathrm{ro}}$,
          and $H_{\mathrm{sp}}$, all backed by domain-separated SHAKE-256 evaluations;
    \item \texttt{threshold\_batch\_ibe}: the five BTD algorithms together with the
          partial-trapdoor and Lagrange-reconstruction logic;
    \item \texttt{tests}: 1\,200 randomised round-trip tests (correctness suite);
    \item \texttt{bench}: wall-clock micro-benchmarks for each algorithm.
\end{itemize}
The source is compiled with \texttt{g++ -std=c++17 -O2} and produces a single
binary that first runs the correctness suite and then, when invoked with an argument,
emits benchmark results.

\subsection{Algorithm mapping}

Table~\ref{tab:algo-map} describes how each step of the formal construction
(Section~\ref{chapter3:Our Threshold Batch Decryption system}) is realised in the prototype
under the toy parameters of Table~\ref{tab:toy-params}.

\begin{table}[h]
\centering
\renewcommand{\arraystretch}{1.25}
\begin{tabular}{|p{0.22\linewidth}|p{0.70\linewidth}|}
\hline
\textbf{Algorithm} & \textbf{Prototype realisation} \\
\hline
$\mathsf{TBIBE.Setup}$ &
  Calls \texttt{ToyHash::trapdoorPublicC} to generate $\mathbf{C}$ (an identity matrix,
  modelling $\mathsf{TrapGen}$). Samples $\mathbf{u}\leftarrow\mathbb{Z}_q^n$ uniformly.
  For each party $j\in[N]$, constructs a partial trapdoor by computing a tail matrix
  $\mathbf{R}$ via \texttt{trapdoorGaussianMatrix} and forming a basis $\mathbf{B}_j$
  that satisfies $\mathbf{C}\mathbf{B}_j\equiv\mathbf{0}$ (modelling $\mathsf{PTrapGen}$).
  \\
\hline
$\mathsf{TBIBE.Enc}$ &
  Samples $\mathbf{s}\leftarrow\mathbb{Z}_q^n$, noise terms $e_1\leftarrow\mathcal{D}_{\chi}$,
  $\mathbf{e}_2\leftarrow\mathcal{D}_{\chi}^{2mk}$, $\mathbf{e}_3\leftarrow\mathcal{D}_{\chi'}^{t+\overline{m}}$.
  Computes $\mathbf{F}_r$ via \texttt{ToyHash::identityMatrixF} (modelling $\mathsf{ExpandLocal}$
  followed by the concatenation with $H_{\mathrm{ro}}(r)$).
  Outputs $(\mathsf{ct}[1],\mathsf{ct}[2],\mathsf{ct}[3])$ as in the formal scheme.
  \\
\hline
$\mathsf{TBIBE.PreDec}$ &
  Derives the batch target $\mathbf{h}=(\mathbf{u},\ldots,\mathbf{u})$ via
  \texttt{ToyHash::batchTarget} (modelling $\mathsf{Expand}$ and the explainable
  $\mathsf{SampleLeft}$ call). Computes the partial preimage $\mathbf{sbk}_j$ via
  \texttt{PartialTrapdoor::samplePreimage} scaled by the Lagrange coefficient
  $\lambda_j(0)$ (modelling $\mathsf{PSampPre}$).
  \\
\hline
$\mathsf{TBIBE.Combine}$ &
  Adds the partial preimages coordinate-wise via $\mathsf{Rec}$:
  $\mathbf{sbk}=\sum_{j\in T}\mathbf{sbk}_j$.
  \\
\hline
$\mathsf{TBIBE.Dec}$ &
  Re-derives $(\mathbf{v}_1,\ldots,\mathbf{v}_\ell)$ via \texttt{identityPreimage}.
  For each ciphertext $i$ computes
  $\mathsf{ct}_i[1]-\mathbf{sbk}^\top\mathsf{ct}_i[2]-\mathbf{v}_i^\top\mathsf{ct}_i[3]$
  and rounds to $\{0,1\}$ by checking proximity to $\lfloor q/2\rfloor$.
  \\
\hline
\end{tabular}
\caption{Mapping of formal BTD algorithms to prototype implementations.}
\label{tab:algo-map}
\end{table}

\subsection{Toy parameters and simplifications}

The prototype is deliberately instantiated at the toy parameter set of
Table~\ref{tab:toy-params} scaled down further to fit a single-threaded
correctness and benchmarking run:
\[
    Q=12289,\quad n=4,\quad k=3,\quad N=5,\quad \ell\in\{1,\ldots,40\}.
\]
Several theoretic components are modelled by deterministic stubs:
\begin{itemize}
    \item $\mathsf{TrapGen}$: replaced by the identity matrix (norm-1 trapdoor);
    \item Gaussian sampler for $\mathbf{T}_{\mathbf{C}_j}$: set to zero (zero-noise
          trapdoor), so the correctness proof holds algebraically without error correction;
    \item $\mathsf{ExpandLocal}/\mathsf{Expand}$: replaced by a deterministic
          hash-derived diagonal matrix $\mathbf{F}_r$ whose structure preserves the
          $\mathbf{C}^{\top}\mathbf{sbk}+\mathbf{F}_r^{\top}\mathbf{v}_i=\mathbf{u}$ invariant.
\end{itemize}
\paragraph{Hash function.}
All three random oracles $H$, $H_{\mathrm{ro}}$, and $H_{\mathrm{sp}}$ are implemented
as domain-separated SHAKE-256 evaluations (Keccak-1600 permutation, FIPS~202).
We vendor the ``more-compact'' C reference from the XKCP project
(CC0 / public domain, $\approx$30 lines): the function
\[
    \mathsf{FIPS202\_SHAKE256}(\textit{in}, \textit{inLen}, \textit{out}, \textit{outLen})
\]
squeezes arbitrarily many output bytes from the Keccak sponge with rate $r=1088$ bits,
suffix byte $\mathtt{0x1F}$, and capacity $c=512$ bits. Each oracle is separated by a
64-bit domain tag prepended to the input before absorption:
\begin{center}
\begin{tabular}{lll}
\hline
\textbf{Oracle} & \textbf{Tag prefix} & \textbf{Output} \\
\hline
$H(r)$ / $H_{\mathrm{ro}}$ row $i$ & $i$ & diagonal scalar in $\mathbb{Z}_q$ \\
$H_{\mathrm{sp}}$ (batch target coord $i$) & $\mathtt{0x0400}\,|\,i$ & element of $\{-1,0,1\}$ \\
\hline
\end{tabular}
\end{center}
The pair $\{H_{\mathrm{ro}}\text{ row }i, H_{\mathrm{sp}}\text{ coord }i\}$ is
distinguished by tag ranges $[0, n)$ versus $[\mathtt{0x400}, \mathtt{0x400}+n)$,
so no two oracles share a tag.
These stubs preserve the \emph{algebraic structure} needed for correctness (the
decryption invariant $\mathbf{C}^\top\mathbf{sbk}+\mathbf{F}_r^\top\mathbf{v}_i
=\mathbf{u}$ holds exactly) while removing statistical Gaussian overhead
that is not yet required for a proof-of-concept evaluation.

\subsection{Correctness test suite}

The correctness suite (\texttt{runRandomizedTests}) executes the following checks:
\begin{enumerate}
    \item Well-formedness of all public-key matrices (entries in $\mathbb{Z}_q$,
          correct dimensions).
    \item Non-degeneracy of noise distributions $\chi,\chi'$ (verified to produce
          non-zero samples within 200 encryptions).
    \item Structural checks on ciphertext and share outputs.
    \item Single-identity correctness: encrypt 0 and 1, combine shares from a
          threshold set, decrypt, assert $m'=m$.
    \item Batch correctness: encrypt a 5-bit message batch, combine, batch-decrypt,
          assert equality.
    \item Input validation: illegal bit values, mismatched sizes, parties outside
          the threshold set, unsorted/duplicate threshold sets, mismatched share
          metadata.
    \item \textbf{1\,000 randomised round-trips}: random seed, random batch size in
          $[1,40]$, random bits, random identities, random threshold set of size $k$;
          assert decryption noise $|\text{noise}|<128$ and $m'=m$ for every ciphertext.
    \item \textbf{200 randomised byte round-trips}: encode arbitrary byte strings
          bit-by-bit, encrypt, decrypt, compare reconstructed bytes.
\end{enumerate}
All 1\,200+ correctness assertions pass deterministically.

\subsection{Experimental setup and results}

Benchmarks were collected on an x86-64 machine (Intel Core i5-11300H \@ 3.10 GHz,
WSL2/Linux, 8 GB RAM) using \texttt{g++ 11.4.0} with \texttt{-O2}.
Timings are wall-clock means over 500 (slow paths) or 2\,000 (fast paths)
repetitions. Sizes are computed from the bit-packed logical representation
($\lceil\log_2 Q\rceil=14$ bits per $\mathbb{Z}_q$ element).

The measured benchmark values are summarised in Table~\ref{tab:eval-metrics} at the
beginning of this chapter.

\subsection{Discussion}

\paragraph{Bottleneck profile.}
With SHAKE-256 as the underlying hash function, the dominant cost across every
algorithm is the Keccak-1600 permutation invoked by the random oracles.
$\mathsf{Enc}$ (71.05 µs) is the most expensive per-message operation: it calls
\texttt{identityMatrixF}, which hashes the identity string once per row
($n=4$ SHAKE-256 calls), then performs two matrix--vector products over
$\mathbb{Z}_q^{4\times 56}$. $\mathsf{PreDec}$ (67.62 µs per party) is
dominated by the $\mathsf{SampleLeft}$ phase (66.79 µs), which absorbs the
packed identity batch into SHAKE-256 once per coordinate of the target vector;
the $\mathsf{PSampPre}$ analogue (0.83 µs) is negligible by comparison.
$\mathsf{Dec}$ (59.35 µs) mirrors $\mathsf{Enc}$ because it must re-derive the
same $\mathbf{v}_i$ vectors via the same SHAKE-256 calls.
$\mathsf{Setup}$ (8.87 µs) and $\mathsf{Combine}$ (0.39 µs) remain cheap
because neither involves per-identity hash evaluations.

\paragraph{Scaling to production parameters.}
The toy PoC uses $Q=12289$, $n=4$, which is orders of magnitude smaller than the
$2^{64}$-secure parameter set of Table~\ref{tab:secure-params} ($\lceil\log_2 q\rceil=127$,
$nk=2560$). Under the production parameters, the matrix dimensions grow to
$\mathbf{C}\in\mathbb{Z}_q^{2560\times 655360}$ (approximately 26.8 GB) and the key vector
$\mathbf{sbk}\in\mathbb{Z}_q^{655360}$ requires $\approx 9$ MB. The dominant cost will be
the $\mathsf{SampleLeft}$ call over a $2560\times (655360+650240)$ matrix together with the
Gaussian sampler operating at width $\sigma_1\approx 2^{45.9}$. Optimising these two
steps---for instance by replacing the Gram--Schmidt orthogonalisation with a parallel
blocked algorithm or exploiting the gadget structure to reduce the effective
dimension---is the most productive direction for improving concrete performance.

\paragraph{What the prototype validates.}
Despite the parameter reductions, the prototype faithfully exercises the full
algebraic path: trapdoor generation $\to$ partial trapdoor distribution $\to$
LWE encryption $\to$ threshold pre-decryption $\to$ Lagrange reconstruction
$\to$ dual-Regev decryption. The 1\,200+ correctness assertions confirm that the
decryption invariant $\mathbf{C}^\top\mathbf{sbk}+\mathbf{F}_r^\top\mathbf{v}_i
=\mathbf{u}$ holds for all tested parameters, providing evidence that the
algebraic design is correct.
